% !TeX program = xelatex
\documentclass[12pt,a4paper]{article}
\usepackage{amsmath,amssymb}
\usepackage{geometry}
\usepackage{booktabs}
\usepackage{graphicx}
\usepackage{xcolor}
\usepackage{tcolorbox}
\usepackage{enumitem}
\usepackage[UTF8]{ctex}
\geometry{margin=2.5cm}

% ============ Custom boxes ============
\newtcolorbox{theorybox}[1][]{
  colback=blue!5!white, colframe=blue!60!black, fonttitle=\bfseries,
  title=#1, sharp corners=south, boxrule=0.8pt
}
\newtcolorbox{keypoint}[1][]{
  colback=red!5!white, colframe=red!60!black, fonttitle=\bfseries,
  title=#1, sharp corners=south, boxrule=0.8pt
}

\title{\textbf{稳定性分析的数学方法详解}\\[5pt]
\large ——以``沪尚回收''回收体系模型为例}
\author{}
\date{}

\begin{document}
\maketitle

% ===================================================================
\section{我们要解决什么问题？}
% ===================================================================

\textbf{问题背景}：上海市``沪尚回收''体系是一个覆盖全市的绿色回收网络，
拥有超过1.5万个基础回收服务点、4300余台智能回收箱。
我们用微分方程描述回收量随时间的变化规律。

\textbf{问题3的核心问题}：
\begin{center}
\fbox{%
\begin{minipage}{0.88\textwidth}
这个回收体系会不会出现``断崖式崩溃''？\\
在怎样的参数范围内，系统是安全的？\\
如何划定安全边界来预警风险？
\end{minipage}}
\end{center}

\textbf{本文要讲清楚的事}：
\begin{enumerate}[label=(\arabic*)]
  \item 这个微分方程模型是怎么来的，每个符号是什么意思
  \item 什么是``平衡点''——系统会停在哪些状态
  \item 什么是``稳定性''——系统受到干扰后，能回到原来的状态吗
  \item 如何用数学严格判断稳定性
  \item 什么是``结构稳定性''——参数变化会导致系统行为发生质的改变吗
  \item 如何划定实用的安全边界
\end{enumerate}

% ===================================================================
\section{微分方程模型：从现实到数学}
% ===================================================================

\subsection{第一步：描述变化率}

我们关心的核心量是 $x(t)$——第 $t$ 天上海市每天的可回收物回收量（吨/天）。

微分方程的基本思想很简单：\textbf{不直接写 $x(t)$ 的表达式，而是写 $x$ 的变化率 $\frac{dx}{dt}$ 与 $x$ 本身之间的关系}。

\medskip
\textbf{类比理解}：银行账户。如果你不直接知道账户余额随时间的变化公式，但你知道``余额越多，利息越多；利息 = 利率 $\times$ 余额''，那么这个关系就是
\[
\frac{dB}{dt} = r \cdot B
\]
其中 $B$ 是余额，$r$ 是利率。这就是一个微分方程——它描述了余额的变化规律，但本身不直接给出余额的值。解这个方程，我们可以得到 $B(t) = B_0 e^{rt}$（指数增长）。

\subsection{第二步：Logistic 方程——自然增长 + 饱和约束}

纯粹指数增长不现实——回收量不可能无限增长，城市的人口和消费能力有限。

最经典的修正来自比利时数学家 Pierre Fran\c{c}ois Verhulst（1838年），他提出了\textbf{Logistic 方程}：
\[
\frac{dx}{dt} = \gamma \cdot x \cdot \left(1 - \frac{x}{K}\right)
\]

逐项解释：
\begin{itemize}
  \item $\gamma > 0$：\textbf{固有增长率}。当 $x$ 远小于 $K$ 时（$x/K \approx 0$），方程近似为 $\frac{dx}{dt} \approx \gamma x$——是指数增长。$\gamma$ 反映了居民基础参与度和政策推动效率。题目给出 $\gamma = 0.022$ /天。
  \item $K > 0$：\textbf{饱和容量}（或称环境承载力）。当 $x = K$ 时，括号内的 $(1 - x/K) = 0$，变化率为零——系统达到上限，不再增长。题目给出 $K = 8011$ 吨/天。
  \item $(1 - x/K)$：\textbf{剩余空间因子}。它像是一个``刹车''——$x$ 越接近 $K$，刹车踩得越重。
\end{itemize}

\subsection{第三步：引入促进和抑制——Q2 的对称拮抗模型}

现实比纯 Logistic 更复杂：
\begin{itemize}
  \item \textbf{促进效应}：站点覆盖率提升、宣传力度加大、投递便捷性改善 → 应该扩大有效容量。
  \item \textbf{抑制效应}：清运不及时、智能回收箱满溢、老年群体操作困难 → 应该降低有效增长速度。
\end{itemize}

Q2 引入两个\textbf{非负}系数来刻画这两类效应：
\begin{itemize}
  \item $\alpha \geq 0$：\textbf{促进系数}。作用于容量——$K_{\text{eff}} = K(1+\alpha)$，促进越强，有效容量越大。
  \item $\beta \geq 0$：\textbf{抑制系数}。作用于增长率——$r = \gamma/(1+\beta)$，抑制越强，有效增长率越低。
\end{itemize}

得到的模型是 Q2 的核心成果：
\begin{equation}\label{eq:model}
\boxed{\frac{dx}{dt} = \frac{\gamma}{1+\beta} \cdot x \cdot \left(1 - \frac{x}{K(1+\alpha)}\right)}
\end{equation}

为书写方便，定义两个\textbf{有效参数}：
\begin{equation}\label{eq:effective}
r = \frac{\gamma}{1+\beta}, \qquad K_{\text{eff}} = K(1+\alpha)
\end{equation}

则方程化为标准 Logistic 形式：
\begin{equation}\label{eq:standard}
\boxed{\frac{dx}{dt} = r \cdot x \cdot \Bigl(1 - \frac{x}{K_{\text{eff}}}\Bigr)}
\end{equation}

\begin{keypoint}[关键约束]
\textbf{$\alpha \geq 0$ 和 $\beta \geq 0$ 不是随意加的——它们由定义本身决定}：
\begin{itemize}
  \item 促进效应用 $\alpha$ 衡量，如果 $\alpha < 0$，那就是``负促进''，这在物理上意味着站点越多、宣传越好反而回收越少，不合理。所以定义上 $\alpha \geq 0$。
  \item 抑制效应用 $\beta$ 衡量，如果 $\beta < 0$，那就是``负抑制''（即促进），这在物理上意味着清运越差反而回收越快，也不合理。所以定义上 $\beta \geq 0$。
\end{itemize}
\textbf{这个看似平凡的事实，将成为整个稳定性分析的关键。}
\end{keypoint}

\subsection*{参数数值一览}

Q2 通过拟合实测数据得到（$\gamma = 0.022, K = 8011$ 是题目给定的固定值）：

\begin{center}
\begin{tabular}{lclc}
\toprule
参数 & 符号 & 数值 & 单位 \\
\midrule
固有增长率 & $\gamma$ & 0.022 & /天 \\
饱和容量 & $K$ & 8011 & 吨/天 \\
\textbf{拟合：有效增长率} & $r$ & $\mathbf{0.008485}$ & /天 \\
\textbf{拟合：有效容量} & $K_{\text{eff}}$ & $\mathbf{8142.22}$ & 吨/天 \\
\midrule
\textbf{导出：促进系数} & $\alpha = K_{\text{eff}}/K - 1$ & $\mathbf{0.0164}$ & 无量纲 \\
\textbf{导出：抑制系数} & $\beta = \gamma/r - 1$ & $\mathbf{1.5929}$ & 无量纲 \\
\bottomrule
\end{tabular}
\end{center}

\textbf{数值含义}：
\begin{itemize}
  \item $r/\gamma = 1/(1+\beta) = 0.386$：抑制效应使有效增长率降至固有增长率的\textbf{仅 38.6\%}——抑制占据绝对主导。
  \item $K_{\text{eff}}/K = 1+\alpha = 1.0164$：促进效应仅将有效容量提升了\textbf{1.64\%}——促进效应较弱。
\end{itemize}

% ===================================================================
\section{预备知识：稳定性理论}
% ===================================================================

在分析模型之前，先把稳定性理论完整地建立起来。
\textbf{每一步都有推导。}

\subsection{什么是``自治微分方程''？}

如果微分方程的右边不显含时间 $t$，只依赖于 $x$ 本身，就称为\textbf{自治微分方程}（autonomous ODE）：
\[
\frac{dx}{dt} = f(x)
\]
其中 $f(x)$ 是某个已知函数。我们的模型正是这个形式：
\[
f(x) = r \cdot x \cdot \left(1 - \frac{x}{K_{\text{eff}}}\right)
\]

\textbf{物理直觉}：自治意味着``游戏规则不随时间改变''——今天和明天的变化规律是一样的。

\subsection{什么是``平衡点''？}

\begin{theorybox}[定义：平衡点（Equilibrium Point）]
如果某个 $x^*$ 满足 $f(x^*) = 0$，则称 $x^*$ 为方程 $\frac{dx}{dt} = f(x)$ 的\textbf{平衡点}（也称不动点、临界点、稳态）。
\end{theorybox}

\textbf{为什么叫平衡点？} 因为如果在 $t=0$ 时系统恰好在 $x^*$，那么 $\frac{dx}{dt}\big|_{x=x^*} = f(x^*) = 0$，
变化率为零——系统永远停在那里不动。

\begin{theorybox}[寻找平衡点的方法]
解代数方程：$f(x) = 0$。这是代数问题，不是微分方程问题——因为平衡点上 $dx/dt = 0$，微分方程退化成了代数方程。
\end{theorybox}

\subsection{稳定性：核心概念}

找到平衡点只是第一步。真正的问题是：\textbf{如果系统没恰好停在平衡点，而是在附近，会发生什么？}

\begin{theorybox}[直观定义：稳定性]
考虑平衡点 $x^*$。假设在 $t = 0$ 时，系统不在 $x^*$，而在 $x^* + \delta$（$\delta$ 是一个微小偏差）：
\begin{itemize}
  \item 如果 $\delta$ 随时间\textbf{衰减到零}（系统回到 $x^*$），则 $x^*$ 是\textbf{渐近稳定的}。
  \item 如果 $\delta$ 随时间\textbf{增长}（系统远离 $x^*$），则 $x^*$ 是\textbf{不稳定的}。
  \item 如果 $\delta$ 既不增长也不衰减（保持常值或振荡），则需要更精细的分析。
\end{itemize}
\end{theorybox}

\textbf{类比}：
\begin{itemize}
  \item 碗底的弹珠 → 稳定平衡（推一下，它会滚回碗底）。
  \item 山顶的弹珠 → 不稳定平衡（推一下，它会滚下山）。
  \item 桌面上的弹珠 → 中性稳定（推一下，它停在新的位置，不回原处也不继续滚）。
\end{itemize}

\subsection{线性化方法}

\subsubsection*{关键思想：在平衡点附近，用切线代替曲线}

考虑 $f(x)$ 在 $x^*$ 附近的\textbf{一阶 Taylor 展开}：
\[
f(x) = f(x^*) + f'(x^*)(x - x^*) + \underbrace{\frac{f''(\xi)}{2}(x - x^*)^2}_{\text{高阶小量，当 } x \to x^* \text{ 时可忽略}}
\]

因为 $x^*$ 是平衡点，$f(x^*) = 0$，所以：
\[
f(x) \approx f'(x^*) \cdot (x - x^*)
\]

\textbf{这就是线性化}：用直线（切线）来近似曲线，只保留一阶项。当 $x$ 离 $x^*$ 足够近时，这个近似是有效的。

\subsubsection*{得到线性微分方程}

令 $u(t) = x(t) - x^*$ 表示偏离平衡点的``偏差量''。则：
\[
\frac{du}{dt} = \frac{dx}{dt} = f(x) \approx f'(x^*) \cdot (x - x^*) = f'(x^*) \cdot u
\]

记 $\lambda = f'(x^*)$，我们得到一个极其简单的微分方程：
\[
\boxed{\frac{du}{dt} = \lambda \cdot u}
\]

\subsubsection*{解这个线性方程}

用\textbf{分离变量法}求解：
\begin{align*}
\frac{du}{dt} &= \lambda u \\
\frac{du}{u} &= \lambda \, dt \quad \text{（两边分离变量）}\\
\int \frac{du}{u} &= \int \lambda \, dt \\
\ln|u| &= \lambda t + C \\
|u| &= e^{\lambda t + C} = e^C \cdot e^{\lambda t}
\end{align*}

令 $u_0 = u(0)$（初始偏差），则 $e^C = |u_0|$，去掉绝对值号后：
\[
\boxed{u(t) = u_0 \cdot e^{\lambda t}}
\]

\subsubsection*{稳定性判据}

\begin{keypoint}[线性化稳定性判据（一维系统）]
对于自治方程 $\frac{dx}{dt} = f(x)$ 和平衡点 $x^*$（满足 $f(x^*) = 0$），
计算 $\lambda = f'(x^*)$：
\[
\begin{cases}
\lambda < 0: & \text{渐近稳定——偏差以 } e^{\lambda t} \text{ 指数衰减到零} \\
\lambda > 0: & \text{不稳定——偏差以 } e^{\lambda t} \text{ 指数放大} \\
\lambda = 0: & \text{线性化无法判断，需要更高阶分析}
\end{cases}
\]
\end{keypoint}

\textbf{为什么 $\lambda$ 叫``特征值''？} 在一维系统中，$\lambda = f'(x^*)$ 就是 Jacobi 矩阵（此时退化为一个数）的特征值。在多维系统中（如捕食者-猎物模型），$f'(x^*)$ 是一个矩阵，我们需要计算其特征值。但这里是一维系统，所以``特征值''就是 $f'(x^*)$ 这一个数字——不需要任何矩阵知识。

\begin{theorybox}[为什么线性化是可靠的？（Hartman--Grobman 定理的通俗表述）]
如果 $\lambda = f'(x^*) \neq 0$（即平衡点是\textbf{双曲的}），那么非线性系统
$\frac{dx}{dt} = f(x)$ 在 $x^*$ 附近的定性行为，与其线性近似
$\frac{du}{dt} = f'(x^*) u$ 完全相同。也就是说：\textbf{只要 $f'(x^*) \neq 0$，
线性化给出的稳定/不稳定结论就是严格正确的}。
\end{theorybox}

\subsection{一维相线：最直观的理解方式}

对于一维自治系统，有一个非常直观的工具——\textbf{相线（Phase Line）}。

\begin{theorybox}[相线的画法]
\begin{enumerate}[label=(\arabic*)]
  \item 在数轴上标出所有平衡点（$f(x) = 0$ 的解）。
  \item 平衡点将数轴分割为若干区间。
  \item 在每个区间内，判断 $f(x)$ 的符号：
  \begin{itemize}
    \item $f(x) > 0$（变化率为正）$\Rightarrow$ $x$ 向\textbf{右}（增大方向）运动 → 画右箭头
    \item $f(x) < 0$（变化率为负）$\Rightarrow$ $x$ 向\textbf{左}（减小方向）运动 → 画左箭头
  \end{itemize}
  \item 观察箭头走向：
  \begin{itemize}
    \item 两边箭头\textbf{都指向}平衡点 → \textbf{稳定}（汇，sink）
    \item 两边箭头\textbf{都背离}平衡点 → \textbf{不稳定}（源，source）
    \item 一边指向、一边背离 → \textbf{半稳定}（鞍点退化）
  \end{itemize}
\end{enumerate}
\end{theorybox}

\textbf{为什么相线与线性化等价？} 因为 $f'(x^*) > 0$ 意味着 $f(x)$ 在 $x^*$ 处是递增的——左边 $f(x) < 0$（箭头向左），右边 $f(x) > 0$（箭头向右）——所以箭头背离平衡点，对应不稳定。$f'(x^*) < 0$ 同理对应稳定。两种方法给出的结论完全一致，但相线更直观，线性化更定量。

% ===================================================================
\section{模型分析第一步：找平衡点}
% ===================================================================

现在回到我们的模型：
\[
f(x) = r \cdot x \cdot \left(1 - \frac{x}{K_{\text{eff}}}\right), \qquad r = \frac{\gamma}{1+\beta}, \quad K_{\text{eff}} = K(1+\alpha)
\]

\subsection{求解 $f(x) = 0$}

\begin{align*}
r \cdot x \cdot \left(1 - \frac{x}{K_{\text{eff}}}\right) &= 0
\end{align*}

因为 $r > 0$（$\gamma > 0$，$1+\beta \geq 1$，所以 $r$ 恒为正），乘积为零的条件是：
\[
x = 0 \quad \text{或} \quad 1 - \frac{x}{K_{\text{eff}}} = 0 \;\Rightarrow\; x = K_{\text{eff}}
\]

代数方程 $f(x) = 0$ 有两个数学解，但 $x$ 的实际定义域为 $[x_0, +\infty)$（回收量从 $x_0 = 6314$ 出发，$f(x) > 0$ 保证 $x$ 单调递增，不会回落），$x_1^* = 0$ 远在定义域之外。\\
\textbf{物理上只存在一个平衡点：}
\begin{equation}\label{eq:equilibria}
\boxed{x^* = K_{\text{eff}} = K(1+\alpha)}
\end{equation}

\textbf{物理含义}：回收量达到有效饱和容量的状态。系统唯一可能停留的位置。

\subsection*{数值代入}
\[
x_1^* = 0, \qquad x_2^* = 8142.22 \text{ 吨/天}
\]

% ===================================================================
\section{模型分析第二步：判断稳定性——线性化方法}
% ===================================================================

\subsection{计算导数 $f'(x)$}

\begin{align*}
f(x) &= r x - \frac{r}{K_{\text{eff}}} x^2 \\[5pt]
f'(x) &= r - \frac{2r}{K_{\text{eff}}} x = r\left(1 - \frac{2x}{K_{\text{eff}}}\right)
\end{align*}

推导过程（逐项求导）：
\begin{itemize}
  \item 第一项 $r x$ 对 $x$ 求导：$\frac{d}{dx}(r x) = r$。
  \item 第二项 $-\frac{r}{K_{\text{eff}}} x^2$ 对 $x$ 求导：$-\frac{r}{K_{\text{eff}}} \cdot 2x = -\frac{2r}{K_{\text{eff}}}x$。
  \item 合并：$f'(x) = r - \frac{2r}{K_{\text{eff}}}x = r\left(1 - \frac{2x}{K_{\text{eff}}}\right)$。
\end{itemize}

\subsection{评估物理平衡点 $x^* = K_{\text{eff}}$ 的稳定性（线性化）}

$f'(x) = r - \frac{2r}{K_{\text{eff}}} x = r\left(1 - \frac{2x}{K_{\text{eff}}}\right)$，在 $x^* = K_{\text{eff}}$ 处：

\[
\lambda = f'(K_{\text{eff}}) = r\left(1 - \frac{2 \cdot K_{\text{eff}}}{K_{\text{eff}}}\right) = r(1 - 2) = -r
\]

因为 $r > 0$：
\[
\boxed{\lambda = -r < 0}
\]

\textbf{结论}：$\lambda < 0$ $\Rightarrow$ $x^* = K_{\text{eff}}$ 是\textbf{渐近稳定的}。

\textbf{物理含义}：如果回收量因为外部冲击偏离了 $K_{\text{eff}}$（偏高或偏低），系统会自动调节、回到 $K_{\text{eff}}$——这是一个``自愈''机制。

\medskip
{\small （附注：数学上 $f(x)=0$ 还有一根 $x=0$，但它远在 $x$ 的定义域 $[x_0, +\infty)$ 之外——回收量从 $x_0 = 6314$ 出发且单调递增，永远不会到达零。$f'(0)=r>0$ 意味着零是不稳定的，这恰好解释了为什么体系不会自发``死掉''，但 $x=0$ 不是系统的有效平衡点。）}

\subsection*{数值代入}
\[
\lambda = -r = -0.008485 < 0 \quad \text{（渐近稳定）}
\]

\subsection{结果总结}

\begin{center}
\begin{tabular}{lccc}
\toprule
平衡点 & 数值（吨/天） & 特征值 $\lambda = f'(x^*)$ & 稳定性 \\
\midrule
$x^* = K(1+\alpha)$ & $8142.22$ & $\lambda = -r = -0.008485 < 0$ & \textbf{渐近稳定} \\
\bottomrule
\end{tabular}
\end{center}

% ===================================================================
\section{模型分析第三步：相线验证}
% ===================================================================

用相线方法交叉验证线性化的结论。$x$ 的物理定义域为 $[x_0, +\infty)$（回收量从 $x_0 = 6314$ 出发且单调递增），$x_0 < K_{\text{eff}}$，平衡点 $x^* = K_{\text{eff}}$ 将定义域切为两个区间：

\[
\underbrace{\qquad [x_0,\; K_{\text{eff}}) \qquad}_{\text{区间 I}}
\;\; \underbrace{\rule{30pt}{0pt}}_{x^*=K_{\text{eff}}}
\;\; \underbrace{\qquad (K_{\text{eff}},\; +\infty) \qquad}_{\text{区间 II}}
\]

$f(x) = r \cdot x \cdot (1 - x/K_{\text{eff}})$ 是三个因子的乘积，$r > 0$ 恒成立。

\subsection*{逐区间、逐因子判断 $f(x)$ 的符号}

\textbf{区间 I：$x_0 \leq x < K_{\text{eff}}$}

\begin{center}
\begin{tabular}{lcc|l}
\toprule
因子 & 表达式 & 符号 & 理由 \\
\midrule
因子 1 & $r$                              & $+$ & $\gamma>0,\;1+\beta>0$ 恒成立 \\
因子 2 & $x$                              & $+$ & $x \geq x_0 = 6314 > 0$ \\
因子 3 & $1 - x/K_{\text{eff}}$           & $+$ & $x < K_{\text{eff}} \;\Rightarrow\; x/K_{\text{eff}} < 1$ \\
\midrule
\multicolumn{3}{l}{\textbf{乘积 $f(x)$}} & \multicolumn{1}{c}{\boxed{$\;+\;$（正）}} \\
\bottomrule
\end{tabular}
\end{center}

$f(x) > 0$ $\Rightarrow$ 变化率为正 $\Rightarrow$ $x$ 随时间\textbf{增大} $\Rightarrow$ 箭头向\textbf{右} $\longrightarrow$

\medskip
\textbf{区间 II：$x > K_{\text{eff}}$}

\begin{center}
\begin{tabular}{lcc|l}
\toprule
因子 & 表达式 & 符号 & 理由 \\
\midrule
因子 1 & $r$                              & $+$ & $\gamma>0,\;1+\beta>0$ 恒成立 \\
因子 2 & $x$                              & $+$ & $x > K_{\text{eff}} > 0$ \\
因子 3 & $1 - x/K_{\text{eff}}$           & $-$ & $x > K_{\text{eff}} \;\Rightarrow\; x/K_{\text{eff}} > 1$ \\
\midrule
\multicolumn{3}{l}{\textbf{乘积 $f(x)$}} & \multicolumn{1}{c}{\boxed{$\;-\;$（负）}} \\
\bottomrule
\end{tabular}
\end{center}

$f(x) < 0$ $\Rightarrow$ 变化率为负 $\Rightarrow$ $x$ 随时间\textbf{减小} $\Rightarrow$ 箭头向\textbf{左} $\longleftarrow$

\subsection*{画出相线}

\[
\begin{array}{cccccc}
& f(x) > 0 & & f(x) < 0 & \\
& \longrightarrow & & \longleftarrow & \\
\hline
x_0 & & x^* = K_{\text{eff}} & & +\infty \\[6pt]
& & \text{渐近稳定} & &
\end{array}
\]

\subsection*{判读}

在物理定义域内只有一个平衡点 $x^* = K_{\text{eff}}$，两边箭头都指向它：
左边箭头向右 $\rightarrow$、右边箭头向左 $\leftarrow$ $\Rightarrow$ \textbf{渐近稳定}。

\textbf{物理直觉}：
\begin{itemize}
  \item 回收量在 $[x_0, K_{\text{eff}})$ 之间 → 还有剩余空间，系统自动增长
  \item 回收量超过 $K_{\text{eff}}$（受外部冲击）→ 超出承载能力，系统自动回落
  \item 不管从定义域内哪一点出发，最终都稳定在 $K_{\text{eff}}$
\end{itemize}

\textbf{这与线性化结论完全一致。}

% ===================================================================
\section{关键发现：为什么系统不会崩溃？}
% ===================================================================

我们得出了最重要的结论：

\begin{center}
\fbox{%
\begin{minipage}{0.88\textwidth}
\textbf{对任意 $\alpha \geq 0$, $\beta \geq 0$}（即促进和抑制都是正向的），\\
系统\textbf{恒收敛}到 $x^* = K_{\text{eff}} > 0$——一个正的有效容量。
\end{minipage}}
\end{center}

\textbf{证明只需一行}：
\[
\beta \geq 0 \;\Rightarrow\; 1+\beta \geq 1 \;\Rightarrow\; r = \frac{\gamma}{1+\beta} \in (0, \gamma]
\;\Rightarrow\; \lambda = -r < 0 \text{ 恒成立}
\]

\textbf{这意味着}：在整个有效的参数空间内（$\alpha \geq 0, \beta \geq 0$），平衡点 $x^*$ 永远是稳定的，稳定性的符号永远不会翻转。系统不会因为参数变化而发生质的改变。

这就是\textbf{结构稳定性}的含义。

% ===================================================================
\section{结构稳定性：更深入的理解}
% ===================================================================

\subsection{什么情况下稳定性会出问题？（但我们不会遇到）}

为了理解结构稳定性的珍贵，我们来看一个反例：\textbf{如果可以将 $\beta$ 扩展到负值会怎样？}

假设我们不考虑 $\beta \geq 0$ 的约束，允许 $\beta$ 取任意实数。那么当 $\beta = -1$ 时：
\[
r = \frac{\gamma}{1+\beta} = \frac{\gamma}{0} \;\rightarrow\; \text{无定义（奇点）}
\]
当 $\beta < -1$ 时：$1+\beta < 0$，所以 $r < 0$，这会导致 $\lambda = -r > 0$——稳定平衡点变成了不稳定的！
系统会从``趋于 $K_{\text{eff}}$''变为``远离 $K_{\text{eff}}$''——这就是\textbf{分岔（bifurcation）}。

\textbf{但}：$\beta$ 是抑制系数，它的定义决定了 $\beta \geq 0$。我们\textbf{永远不会}接近 $\beta = -1$。

\subsection{结构稳定性的定义}

\begin{theorybox}[结构稳定性（通俗定义）]
一个微分方程系统是\textbf{结构稳定}的，如果对参数（比如 $\alpha, \beta$）的微小改变，
系统的\textbf{定性行为}不发生变化——平衡点的个数不变，每个平衡点的稳定/不稳定性质也不变。
\end{theorybox}

换个说法：结构稳定意味着系统的``拓扑相图''在参数扰动下保持不变。
如果参数的变化会导致平衡点突然出现、消失，或稳定性翻转，
就发生了\textbf{分岔}——系统就不是结构稳定的。

\subsection{我们的模型为什么是结构稳定的？}

因为稳定性由 $\lambda = -r = -\gamma/(1+\beta)$ 的符号决定。
而 $\text{sign}(\lambda) = \text{sign}(-\gamma/(1+\beta))$。

在有效参数空间 $\{(\alpha, \beta) \mid \alpha \geq 0, \beta \geq 0\}$ 中：
\[
\gamma > 0,\; 1+\beta \geq 1 > 0 \;\Rightarrow\; r > 0 \;\Rightarrow\; \lambda < 0
\]

\textbf{符号永远不会翻转。} 这就是结构稳定性的数学本质。

\begin{keypoint}[关键结论]
\begin{center}
\textbf{``沪尚回收''体系在数学上不可能发生断崖式崩溃。}\\[3pt]
只要促进效应和抑制效应保持其设计的物理含义（恒为正数），\\
回收量将始终趋向正平衡值 $K_{\text{eff}}$。
\end{center}
\end{keypoint}

\subsection*{那风险在哪里？}

风险不在于数学崩溃，而在于\textbf{性能退化}：

\begin{itemize}
  \item 如果 $\beta$ 越来越大（运维质量越来越差），$r = \gamma/(1+\beta)$ 会趋近于零——系统理论上仍然收敛到 $K_{\text{eff}}$，但收敛得\textbf{极慢}，实际运营中看起来像``停滞''。
  \item 如果 $\alpha$ 趋近于零，$K_{\text{eff}} \to K$，有效容量缩水到基准值。
  \item 现实中，回收量在``缓慢爬升''时，居民看不到进展、管理者失去信心——这本身就会引发恶性循环。
\end{itemize}

\textbf{所以我们需要实用安全边界——下一节讨论。}

% ===================================================================
\section{实用安全边界：虽然不是分岔，但不能太慢}
% ===================================================================

结构稳定性给了我们定性的安全感，但运营管理还需要定量的警戒线。

\subsection{增长率约束：不能收敛得太慢}

设 $r_{\min}$ 为\textbf{最低可接受的有效增长率}。我们要求：
\[
r = \frac{\gamma}{1+\beta} \geq r_{\min}
\]

\textbf{为什么设这个约束？} 如果 $r$ 太小，系统虽然理论上会收敛到 $K_{\text{eff}}$，
但需要极长时间。例如 $r = 0.001$ 时，恢复到平衡值 $99\%$ 需要约 $\ln(99)/0.001 \approx 4595$ 天（约 12.6 年）
——这在实践中是不可接受的。

解出 $\beta$ 的上界：
\begin{align*}
\frac{\gamma}{1+\beta} &\geq r_{\min} \\[3pt]
\gamma &\geq r_{\min}(1+\beta) \\[3pt]
\frac{\gamma}{r_{\min}} &\geq 1+\beta \\[3pt]
\beta &\leq \frac{\gamma}{r_{\min}} - 1
\end{align*}

取 $r_{\min} = 0.005$（基于运营管理的最低要求：\\
大约需要 $-\ln(0.01)/0.005 \approx 921$ 天恢复到 $99\%$ 的饱和水平）：
\[
\boxed{\beta_{\max} = \frac{0.022}{0.005} - 1 = 3.40}
\]

当前 $\beta = 1.59$，安全裕度 $= 3.40 - 1.59 = 1.81$——\textbf{裕度充足}。

\subsection{容量约束：不能低于基准}

设 $K_{\min}$ 为最低可接受的有效容量。我们取 $K_{\min} = K$（至少不能比纯自然的饱和值还低）：
\[
K_{\text{eff}} = K(1+\alpha) \geq K \;\Rightarrow\; 1+\alpha \geq 1 \;\Rightarrow\; \alpha \geq 0
\]

\textbf{这个约束由 $\alpha$ 的定义自然满足}——$\alpha \geq 0$ 就是促进系数的定义本身。

\subsection{安全运行范围总表}

\begin{center}
\begin{tabular}{lcccc}
\toprule
参数 & 下界 & 上界 & 当前值 & 状态 \\
\midrule
$\alpha$（促进） & $0$（定义） & — 无上界 & $0.0164$ & 安全 \\
$\beta$（抑制） & $0$（定义） & $3.40$（$r_{\min}$ 约束） & $1.59$ & 安全，裕度 $1.81$ \\
\bottomrule
\end{tabular}
\end{center}

% ===================================================================
\section{解耦调控：模型的一个优雅特性}
% ===================================================================

观察有效参数的定义：
\[
r = \frac{\gamma}{1+\beta}, \qquad K_{\text{eff}} = K(1+\alpha)
\]

\textbf{发现}：
\begin{itemize}
  \item $\alpha$ 只出现在 $K_{\text{eff}}$ 中，不影响 $r$。
  \item $\beta$ 只出现在 $r$ 中，不影响 $K_{\text{eff}}$。
\end{itemize}

\textbf{这意味着决策者可以独立调控两个维度：}
\begin{itemize}
  \item \textbf{想提升容量} → 增大 $\alpha$（扩大站点覆盖、改善投递便捷性）→ $K_{\text{eff}}$ 增大，但 $r$ 不变。
  \item \textbf{想加快收敛速度} → 减小 $\beta$（改善清运效率、降低运维摩擦）→ $r$ 增大，但 $K_{\text{eff}}$ 不变。
\end{itemize}

这就像汽车的两个独立踏板——油门（$\alpha$）控制最终能跑多快（容量），
离合器（$\beta$）控制换挡的平顺程度（收敛速度）。
两个踏板互不干扰，可以分别优化。

% ===================================================================
\section{理论总结：完整逻辑链}
% ===================================================================

把从建模到最终结论的逻辑链串起来：

\begin{enumerate}[label=\textbf{步骤 \arabic*}.]
  \item \textbf{建立微分方程}：
        \[
        \frac{dx}{dt} = r \cdot x \cdot \left(1 - \frac{x}{K_{\text{eff}}}\right), \quad r = \frac{\gamma}{1+\beta}, \; K_{\text{eff}} = K(1+\alpha)
        \]
        其中 $\alpha \geq 0$（促进系数），$\beta \geq 0$（抑制系数），由定义保证。

  \item \textbf{找平衡点}：令 $f(x) = r x (1 - x/K_{\text{eff}}) = 0$。$x$ 的定义域为 $[x_0, +\infty)$，
        只有 $x^* = K_{\text{eff}}$ 在定义域内。（数学上另有一根 $x=0$，远在定义域外。）

  \item \textbf{线性化}：计算 $f'(x) = r(1 - 2x/K_{\text{eff}})$，在 $x^* = K_{\text{eff}}$ 处：
        \[
        \lambda = f'(K_{\text{eff}}) = -r < 0 \;\Rightarrow\; x^* \text{ 渐近稳定}
        \]
  \item \textbf{结构稳定性论证}：因为 $\beta \geq 0 \Rightarrow r = \gamma/(1+\beta) > 0$ 恒成立，
        所以 $\lambda = -r < 0$ 在整个有效参数空间内恒成立。
        \textbf{不存在任何参数取值使稳定性翻转}——系统是结构稳定的。

  \item \textbf{实用安全边界}：虽然不会崩溃，但运营效率要求 $r \geq r_{\min}$，
        推导出 $\beta \leq \gamma/r_{\min} - 1 = 3.40$（取 $r_{\min} = 0.005$）。

  \item \textbf{解耦调控}：$\alpha$ 只控容量、$\beta$ 只控速度，可独立优化。
\end{enumerate}

\begin{keypoint}[核心方法总结]
一维自治系统稳定性分析的核心工具只有三个：
\begin{itemize}
  \item \textbf{求解 $f(x) = 0$}（代数求解）——找平衡点
  \item \textbf{求导 $f'(x)$}（微积分求导）——线性化
  \item \textbf{判断导数的符号}（正负号分析）——稳定性判据
\end{itemize}
这就是一维自治系统稳定性分析的全部。多维系统才需要矩阵特征值；
一维系统只需要一个导数和一个符号判断。
\end{keypoint}

\end{document}
